\section*{ID9234: Relation: δω/ω105δω/ω∼10Θ−5δω/ω105aspredicted.}
\paragraph{Metadata.}
PiID: 5765615144632 | RTEID: RTE:P35857.6727:T2.5337 | UUID: 60926b15-cbf8-5f6d-91f6-1ce7274044a5

\paragraph{Canonical Statement.}
al pattern by leveraging helicity conservation and KK mode relationships. Dependencies: IDs 314–316. Derivable Consequences: (1) Allows determination of compactification radius RKKR\_\{\textbackslash{}text\{KK\}\}RKK from spectral data. (2) Enables inversion of hidden extradimensional excitations. (3) Provides a basis for experimental tests of reverse harmonic protocols. Falsifiability / Test Handle: Apply the algorithm to simulated spectra with known initial fields; verify accurate reconstruction and sensitivity \$δω/ω105\textbackslash{}delta\textbackslash{}omega/\textbackslash{}omega\textbackslash{}sim10Θ\{−5\}δω/ω105aspredicted.\$ 2.0.169 ID 318: Constrained Dipole Field Configuration Type: Definition Formal Statement: Two magnetic dipoles of equal magnitude \$µ0\textbackslash{}mu\_0µ0,\$ oriented along the zzzaxis but with opposite direction, are located at r1=(0,0,d/2)r\_1=(0,0,-d/2)r1=(0,0,d/2) and r2=(0,0,d/2)r\_2=(0,0,d/2)r2=(0,0,d/2). The field produced by a single dipole at position r0r\_0r0 is \$Bdip(r;µ,r0)=(µ0/4π)[3(rr0)(µ(rr0))/rr0

\paragraph{Canonical Equation.}
\[
δω/ω105\delta\omega/\omega\sim10Θ{−5}δω/ω105aspredicted.
\]

\paragraph{Dictionary.}
Relation: δω/ω105δω/ω∼10Θ−5δω/ω105aspredicted. — δω/ω105δω/ω∼10Θ−5δω/ω105aspredicted.

\paragraph{Encyclopedia.}
Relation: δω/ω105δω/ω∼10Θ−5δω/ω105aspredicted. is a scaling / order-of-magnitude relation, written δω/ω105δω/ω∼10Θ−5δω/ω105aspredicted.. It expresses a scaling or proportionality, capturing how the quantity grows with its parameters. It is built from δ, ω, a, s, p, r. Within the Trawin Zero Point registry it serves as one element of the framework's mathematical narrative.
